\section{Práctica de ingeniería y ajuste de hiperparámetros}

\subsection{Replay buffer y estrategias de muestreo de datos}

Los modelos Actor-Critic off-policy modernos suelen usar un replay buffer con prioridad, frame stacking y multi-environment rollout. Las notas de clases anteriores también han insistido reiteradamente en que la estrategia de muestreo de datos determina directamente la señal de entrenamiento del critic.

\begin{knowledgebox}{Pequeños trucos del Prioritized replay}
Se usa el valor absoluto del TD error como priority, lo que aumenta la probabilidad durante el sampling, y se multiplica el loss por el importance weight correction para garantizar que el update no favorezca los ejemplos con error alto.
\end{knowledgebox}

\subsection{Práctica de Batching y del retorno de múltiples pasos}

El tamaño de lote, la frecuencia de muestreo y el n-step, en combinación, constituyen una dirección básica del ajuste de hiperparámetros. La siguiente tabla resume las combinaciones más comunes.

\begin{table}[H]
\centering
\begin{tabular}{p{0.25\textwidth}p{0.35\textwidth}p{0.28\textwidth}}
\toprule
Elemento de configuración & Valores comunes & Efecto en el entrenamiento \\
\midrule
Batch size & 256--4096 transitions & Cuanto mayor, menor el variance, pero el rendimiento del hardware es el cuello de botella \\
n-step & 1--5 & Con más pasos el sesgo es menor, pero aparece drift; se usa comúnmente un $\lambda$ mix \\
Update frequency & 1 update per 1--4 steps & Menos update reduce el off-policy bias \\
Replay ratio & 1--4 & >1 significa que se actualiza varias veces después de cada muestreo; el learning rate de cada paso debe reducirse en la misma proporción \\
\bottomrule
\end{tabular}
\caption{Combinaciones de hiperparámetros que se ajustan con frecuencia en la práctica}
\end{table}

\subsection{Ajuste de hiperparámetros y diagnóstico de fallas}

\begin{figure}[H]
\centering
\includegraphics[width=0.82\textwidth]{slide-15.jpg}
\caption{Diagrama de flujo de entrenamiento mostrado en la Lecture 4: varios worker extraen datos, el critic ajusta el objetivo y el Actor actualiza la política.}
\end{figure}

Cuando el entrenamiento presenta divergence, primero se observan las three signals: critic loss, KL y reward curve. Además, los failure mode comunes incluyen:
\begin{itemize}
    \item Critic overestimate: la regresión falla y hace que el advantage sea siempre positivo;
    \item Replay buffer stale: el entorno cambia pero los datos del buffer siguen siendo los mismos;
    \item Exploration collapse: el entropy cae rápidamente a 0.
\end{itemize}

\begin{warningbox}{Procedimiento rápido de reinicio cuando aparece divergence}
Primero regrese la política al checkpoint anterior, desactive el entropy decay, reduzca el learning rate y luego recupérela gradualmente. No aumente el volumen de datos de una sola vez, o se acelerará el desastre.
\end{warningbox}

\subsection{Resumen de la sección}

En la práctica, el muestreo de datos, el batching y la programación de hiperparámetros constituyen el problema sistemático de la estabilidad del entrenamiento de Actor-Critic. Monitorear el replay ratio, el entropy y el critic loss son los datos de primera mano para el ajuste de hiperparámetros.
